\documentclass[11pt]{article}
\usepackage{amsmath,amssymb,mathtools}
\usepackage{fontspec}
\usepackage{unicode-math}
\setmainfont{Latin Modern Roman}
\setmathfont{Latin Modern Math}
\usepackage{xurl}
\usepackage[unicode,colorlinks=true,linkcolor=blue,urlcolor=blue,pdftitle={Problem statement}]{hyperref}
\newcommand{\sref}[2]{\hyperlink{source:#1:#2}{[S#2]}}
\newcommand{\cataloguescope}[1]{\paragraph{Mathematical statement.}#1}
\begin{document}
\section{Research Problems in Function Theory — Problem 2.19}\label{statement}

\subsection{Definitions and mathematical statement}

\paragraph{Shared notation.}$\mathbb N=\{1,2,\ldots\}$, and $\mathbb Z,\mathbb Q,\mathbb R,\mathbb C$ denote the integers, rationals, reals and complex numbers. Any other conventions are specified locally.


Fix $A,B>0$ and $K\ge0$. Let $\mathcal E(A,B,K)$ be the class of entire functions for which there is a constant $C_f>0$ satisfying
\[
 |f(z)|\le C_f e^{K|z|}\quad(z\in\mathbb C),\qquad
 |f(x)|\le A\quad(x<0),\qquad |f(x)|\le B\quad(x>0).
\]
Thus $K$ bounds the exponential type, while the multiplicative constant in the global growth bound may depend on $f$.

\paragraph{Statement recorded in the supplied source.}
Find the sharp pointwise bound
\[\Phi_{A,B,K}(z)=\sup\{|f(z)|:f\in\mathcal E(A,B,K)\}\]
in terms of $A,B,K$ and $z$. \sref{statement}{2}
For equal half-line bounds the familiar bound is $|f(x+iy)|\le A e^{K|y|}$. The problem asks for the unequal-bound case; Update 2.19 distinguishes a subharmonic upper estimate from the sharp entire-function problem. \sref{statement}{2}
\subsection{Short English statement}
Determine the best possible complex-plane bound for an entire function of bounded exponential type when its bounds on the two real half-lines may differ.
\subsection{Sources}
\begin{description}
\item[\hypertarget{source:statement:1}{S1}] ULAM.AI and the UnsolvedMath Contributors, \emph{UnsolvedMath: A Curated Collection of Open Mathematics Problems}, catalogue entry. CC BY 4.0. \url{https://huggingface.co/datasets/ulamai/UnsolvedMath}.
\item[\hypertarget{source:statement:2}{S2}] W. K. Hayman and E. F. Lingham, \emph{Research Problems in Function Theory} (2018), Chapter 2, Problem 2.19 and Update 2.19. \url{https://arxiv.org/abs/1809.07200}.
\end{description}
\label{end:statement}
\end{document}
